\begin{abstract}
\noindent
Literature-derived datasets are an appealing route to data-driven design of MXene
supercapacitor electrodes, but they are small and heterogeneous. We compile 118 usable
gravimetric-capacitance observations for Ti\textsubscript{3}C\textsubscript{2}T\textsubscript{x}
in H\textsubscript{2}SO\textsubscript{4} from 40 publications, standardise thirteen
structural and processing descriptors with unit-aware parsing, and evaluate prediction under
strict publication-grouped cross-validation with a full target-leakage audit. Cross-paper
transferability is weak
($R^2 = 0.03 \pm 0.12$, MAE
$= 82$\,F\,g\textsuperscript{-1}). A variance decomposition localises
the cause: 64\% of the variance in capacitance and a mean of
87\% of the variance in the descriptors lies \emph{between}
publications, and several descriptors are constant within each publication
(ICC$(1)\approx1$). We then test whether Bayesian-network augmentation
(\textsc{DataSynthesizer}, correlated-attribute mode) relieves the small-sample constraint,
fitting the generator inside each training fold only and sweeping four augmentation ratios,
three network degrees and three generator seeds. It does not: the best condition improves MAE
by 6.0\,F\,g\textsuperscript{-1}, less than the
12.3\,F\,g\textsuperscript{-1} obtained by re-drawing the synthetic set alone, with
a non-monotonic response in the amount of synthetic data. The generator itself was faithful
--- no marginal rejected by a Kolmogorov--Smirnov test, Jensen--Shannon distance
0.08 on categorical frequencies, dependency structure
reproduced to $\approx0.06$, and no memorisation of real rows
--- so the null result reflects the nature of the bottleneck rather than a generator failure:
resampling a learned joint distribution cannot supply information about experimental domains
the source papers never visited. SHAP attribution is nevertheless highly stable under this
resampling (Spearman $\rho = 0.90$, top-5 overlap
5/5), which licenses using the real-only ranking for descriptor
prioritisation. Finally we assess what physics the corpus can enforce. A classical
PDE-based physics-informed neural network is not justified --- potential window, discharge
time, current and any time axis are entirely absent --- but two algebraic capacitance closure
identities hold to a median relative error below 1\%, allowing areal and volumetric
capacitance to serve as auxiliary supervised outputs under a hard consistency constraint. We
therefore recommend a hybrid physics-constrained multi-task regression, and identify the
potential window and discharge time as the highest-value missing measurements.
\end{abstract}

\vspace{4pt}
\noindent\textbf{Keywords:} MXene; supercapacitor; gravimetric capacitance; SHAP; Bayesian
network; synthetic data augmentation; domain shift; physics-guided machine learning.
